\chapter{Related literature}\label{ch:lit}
\chapterlead{The theory of why arbitrage is limited, how the band of a futures basis is estimated, what short-sale
constraints do to prices, and what earlier work found in India.}

\section{Why arbitrage is limited}
Shleifer and Vishny (1997) start from the fact that arbitrage is done by specialists who invest other people's money.
When a gap widens, the arbitrageur loses money and investors withdraw funds, exactly when the opportunity is best; and
specialised arbitrageurs avoid positions whose returns are volatile. De Long, Shleifer, Summers and Waldmann (1990) show
that risk created by noise traders can keep prices from fundamentals, although the result needs that risk to be
market-wide and is not robust to changes in the model (Loewenstein and Willard, 2006). Gromb and Vayanos (2010) bring the
main mechanisms into one model: a cost of short selling adds a wedge only when arbitrageurs must be short, so the band
around fair value is one-sided, and margins make capital a binding constraint. The model of Chapter~\ref{ch:theory}
follows this structure and adds a position limit under which every participant may only reduce positions.

Two classic papers explain why short-sale constraints matter. Miller (1977) argues that when opinions differ and
pessimists cannot sell short, prices reflect the optimists. Diamond and Verrecchia (1987) show that with rational
traders constraints slow the incorporation of private information, especially bad news. Ofek, Richardson and Whitelaw
(2004) find put-call parity violations in the direction short-sale costs predict, and Stambaugh, Yu and Yuan (2015) show
that the low returns of volatile stocks come from the overpriced side, where arbitrage is hardest.

\section{The band and the evidence from India}
Trading costs create a band around the cost-of-carry price within which arbitrage does not pay. Dwyer, Locke and Yu
(1996) estimate it with a three-regime threshold model of S\&P 500 futures and cash; Tse (2001) argues that arbitrageurs
with different costs make the response smooth; Ters and Urban (2020) show that the arbitrage cost itself can be estimated
precisely and propose comparing bands before and after regulatory changes.

Indian studies (Table~\ref{tab:litbands}) estimate intraday bands over short samples. Aggarwal (2015) finds a band of
about $-15$ to $+3$\bp\ that moved with funding liquidity; Shankar, Sankar and Kiran Kumar (2019) find thresholds of $-39$
and $-10$\bp\ and a lower edge that reacts more to volatility, a sign of a constraint on the short side. Saurav, Srivastava
and Agarwalla (2026) report, in their abstract, that the basis turns negative before earnings announcements where
short-sale constraints are tight. In the cross-section, Sehgal, Vasishth and Deisting (2024) find that volatile stocks
earn \emph{higher} returns in India, the opposite of the United States, so results from elsewhere need checking here.

\begin{table}[t]
\caption[Earlier estimates of the futures-cash band in India]{Earlier Indian estimates describe intraday bands; this thesis measures the overnight band}\label{tab:litbands}
\exsub{Studies of single-stock futures against their underlying stocks on the National Stock Exchange}
\small\begin{tabularx}{\linewidth}{@{}K{3.4cm}P{2.9cm}P{2.4cm}L@{}}
\toprule
\thd{Study} & \thd{Sample} & \thd{Data} & \thd{Band} \\
\midrule
Aggarwal (2015) & 141 stocks, Mar to Aug 2013 & One-second order book & About $-15$ to $+3$\bp, per stock and day \\ \rowrule
Shankar, Sankar and Kiran Kumar (2019) & 102 stocks, 2007 to 2009 & Five-minute prices & $-39$ to $-10$\bp, weekly \\ \rowrule
Slivka, Shah and Wu (2011) & Cost calculation & Fee schedules & About twice the width in mature markets \\ \rowrule
This thesis & 360 stocks, 2015 to 2026 & Daily volume-weighted prices & $-34.5$ to $+23.9$\bp\ around each stock's median \\
\bottomrule
\end{tabularx}
\exsource{Source: the studies cited; Chapter~\ref{ch:band} for this thesis.}
\end{table}

\section{What this thesis adds}
To my knowledge no study has estimated the effects of India's F\&O ban. I model it as a reduce-only rule for every
participant, test its predictions with a plan fixed in advance, estimate the band over a long daily panel of stocks
rather than a short intraday sample, relate it to dated changes in costs, capital and position limits, and study
short-sale constraints across the whole cash market through entry into the F\&O segment and a constraint index for every
listed stock.
